\chapter{Implementation}
\label{ch:implementation}

\polyforge{} is a working system, not a simulator with a paper wrapped
around it. This chapter records what was built, the engineering gates it is
held to, the deliberate deviations from the original roadmap, and the wiring
that makes the live evaluation arm honest.

\section{Components and code layout}

\begin{table}[htbp]
\centering\small
\begin{tabularx}{\textwidth}{@{}lX@{}}
\toprule
\textbf{Component} & \textbf{Code and role} \\
\midrule
Control plane & \evfile{cmd/control-plane}, \evfile{internal/}: tenant /
project / API-key APIs, PostgreSQL with row-level security (SQLite dev
path), rate limiting, telemetry ingestion, canary and isolation-ladder admin
APIs. \\
AI gateway & \evfile{cmd/ai-gateway}, \evfile{internal/ai}: multi-backend
chat routing (plan/length/budget), semantic cache with cost-aware eviction,
vector search, agent loop. \\
Classifier & \evfile{internal/classifier}, \evfile{cmd/classifier-train}:
12-feature windows $\to$ 5-class taxonomy, online labels every 10\,s, PSI
drift detection. \\
\jcac{} planner & \evfile{services/planner} $+$
\evfile{research/jcac\_sim}: the service imports the simulator's controller,
so offline and online planning are the same code. \\
Operator & \evfile{cmd/operator}, \evfile{deploy/helm/polyforge-operator}:
four CRDs, guardrailed reconciliation, ordered finalizer teardown,
OTel-traced reconciles. \\
Evaluation & \evfile{eval/}: YAML-driven harness, five tuned baselines,
DuckDB results, validation and spot-check tooling; kind-cluster live
backend. \\
\bottomrule
\end{tabularx}
\caption{Implementation map.}
\label{tab:impl-map}
\end{table}

\section{Engineering gates}
\label{sec:impl-gates}

Every session's work passes the same battery before it is committed:
\texttt{gofmt} clean, \texttt{go vet}, \texttt{go build},
\texttt{go test ./...\ -count=1} across 26 Go packages, the Python suites
(harness, planner, analysis --- 41 tests at the time of writing), and
OpenAPI linting. CI additionally runs the race detector
(\texttt{-race -covermode=atomic}) --- which cannot run on the development
machine --- a coverage gate, golangci-lint, and the PostgreSQL row-level
security suite against a live PostgreSQL service. Two verification norms
follow from that split and are kept throughout the project record: tests
that only CI can execute are reported as \emph{unproven until a green CI
run}, and the session log (\evfile{docs/SESSION\_LOG.md}) records how every
change was verified.

Security-relevant details received the same treatment as research claims.
API-key identifiers are drawn independently of the secret (an audit found
the ID leaked 64 of 192 entropy bits); keys are hashed with SHA-256 rather
than argon2id --- sound because the secrets are 192-bit random strings where
brute force is infeasible, a deliberate deviation recorded as an ADR
precisely because it is defense bait; a panic-recovery middleware returns
structured 500s without leaking stack details; the authentication packages
are held at ${\sim}95\%$-plus coverage.

\section{Deliberate deviations from the roadmap}

The original roadmap named specific tools; the implementation deviates in
three places with equivalent outcomes, each recorded rather than silently
``fixed'': the chi router is replaced by Go 1.22$+$ stdlib
\texttt{http.ServeMux} method patterns; sqlc and golang-migrate are replaced
by hand-written SQL behind repository interfaces plus an idempotent
migrator (type-safe access and versioned schema still hold, RLS proven by
test); argon2id by SHA-256 as above. The evaluation deviates from the
roadmap's p99 in reporting p95 (Section~\ref{sec:method-metrics}).

\section{The live evaluation arm, and two desk-found defects}
\label{sec:impl-live}

The live campaign (Section~\ref{sec:eval-live}) runs on a kind cluster:
harness-built images, a k6 load generator, real HPA for the baseline arm,
and the operator $+$ planner for the \jcac{} arm. Wiring the \jcac{} arm
surfaced two defects at the desk that would each have silently invalidated
the cluster session:

\begin{enumerate}
  \item \textbf{Demand plumbing that could never have authenticated.} The
        operator's feature-demand source sent its token as a bearer, which
        the control plane treats strictly as a JWT; every read would have
        been a 401, and the \jcac{} loop \emph{degrades by design} into
        silent fallback --- the arm would have run ``live'' while planning
        nothing. The features endpoint (that endpoint only) now accepts the
        platform admin key, unit-tested in both directions (the admin key
        reads any tenant's features; it still cannot touch other tenant
        routes).
  \item \textbf{Last-writer-wins actuation on a shared target.} All eight
        eval tenants' Policies name one shared Deployment; each reconcile
        overwrote the others' replica counts, which would have pinned the
        target at a single tenant's share --- \jcac{} crippled by wiring and
        mislabeled as \jcac{}. Policies resolving to the same target now sum
        their clamped contributions, with deletion re-enqueuing siblings;
        unit-tested.
\end{enumerate}

Both fixes carry a broader honesty mechanism: the harness ends operator
installation with \texttt{kubectl wait -{}-for=condition=Applied} on every
Policy, so a run in which the operator never actuated \emph{fails before
any load is generated} --- the ``the controller was actually in the loop''
gate is executable, not a promise. Capacity parity is a cell property under
test: every arm starts at the same replica count and is ceilinged by the
same cluster limits, so no arm can win by being handed more cluster.

\section{Scale behavior of the deployed planner}
\label{sec:impl-scaling}

The deployed \texttt{PlannerCore.plan} path was driven from 8 to 256 tenants
on one laptop core (\evfile{research/analysis/PLANNER\_SCALING.md}): fitted
growth exponent \textbf{1.45} --- the joint fairness term couples tenants,
and super-linearity is the measured price of jointness --- with p95 cycle
time first exceeding the operator's 3\,s request timeout at 128 tenants and
the 10\,s control period at 256. Past the timeout the failure mode is the
designed one: the operator holds the last good plan and reports
\texttt{fallback} status. Scaling levers past the crossover (planning cells,
incremental fairness partial sums, a faster inner loop) are named with the
measurement. Coordination costs outside the planner --- CR write fan-out,
telemetry aggregation --- are not covered by this microbenchmark.
